\section{Assemble and store a complete animation}
A record encoder owns one image transition. The file encoder must give those
transitions a beginning, order and exact end. Its source is a packed array of
native images; its destination is a byte sequence with one global header followed
by exactly the declared number of records. No path or file descriptor belongs to
this computation. Keeping the staged result available allows the existing reader
to reject a construction error before an output file is opened.

Implement \texttt{video\_encode\_file} with source address in \texttt{\$a0}, exact
source byte count in \texttt{\$a1}, staging address in \texttt{\$a2} and a native
96-byte context in \texttt{\$a3}. The source contains one through 256 consecutive
64 by 64 images, each 16\,384 bytes of native \texttt{0x00RRGGBB} words. Require
four-byte source and eight-byte context alignment. The staging address may be
unaligned. All caller-supplied windows and the private stack are disjoint, except
for the deliberate staging alias described below.

\begin{table}[htbp]
\centering
\caption{The complete-file encoder's native context.}
\label{tab:video-file-encode-context}
\begin{tabular}{r p{1.5in} p{2.8in}}
\toprule
Offset & Field & Meaning \\
\midrule
0 & Staging capacity & 12\,328 through 7\,344\,152 bytes. \\
8 & Frame count & One through 256; source bytes equal this count times 16\,384. \\
16 & Full-frame period & One through 64. \\
24 & Frame duration & One through 1000 milliseconds. \\
32, 40 & Previous address, capacity & Four-byte-aligned working image; at least 16\,384 bytes. \\
48, 56 & XOR address, capacity & Four-byte-aligned scratch; at least 16\,384 bytes. \\
64 & Produced file bytes & Output on complete success only. \\
72 & Produced records & Output on complete success only. \\
80--95 & Retained suffix & Leave unchanged. \\
\bottomrule
\end{tabular}
\end{table}

The capacity ceiling follows the format's maximum admitted record, not this
particular encoder's usual choice: $24+256\times28\,688=7\,344\,152$ bytes.
The minimum admits a global header and one full record. It does not guarantee room
for every declared image. With capacity at least $24+n\times28\,688$, every valid
$n$-frame input must succeed. At a smaller admitted capacity, a valid compact
representation may succeed or a selected representation may fail to fit.
Completed grading accepts equivalent valid full or delta choices, including extra
full records. Every required periodic full record remains mandatory.

Return zero with the exact file count on success, or one with a zero count on
completed failure. Development status four still means an unfinished region.
The source remains immutable. On failure, the whole caller context remains unchanged,
including old output fields. Staging, working previous pixels and XOR scratch may
change within their active windows. Those temporary bytes are not restored by a
failed call. Preserve an old animation in a separate area if it must survive a
replacement attempt. Preserve saved registers and sixteen-byte stack alignment
across nested record calls.

\begin{table}[htbp]
\centering
\caption{Storage roles during complete-file encoding.}
\label{tab:video-file-encode-storage}
\begin{tabular}{p{1.5in} p{3.5in}}
\toprule
Storage & Lifetime and interpretation \\
\midrule
Packed source & Immutable native words; 16\,384 bytes per image. \\
Staging & Mutable serialized header and records; publish only an exact complete extent. \\
Previous and XOR & Temporary native words; their active contents may change on failure. \\
Caller context & Inputs retained; two output fields change only on success. \\
Private record context & New sequence zero for each file; lives in the current stack frame. \\
\bottomrule
\end{tabular}
\end{table}

From \texttt{drafts/video-file-encode-construction}, attach your preceding record
encoder with \texttt{attach.py --source record-work.asm --output file-work.asm}.
The importer retains that working file and gives the existing file arena the alias
\texttt{file\_encode\_output}. It does not reserve a second maximum-sized staging
arena. Both names identify the same storage; an old file and its directory cease
to be usable while new encoding overwrites those bytes. The packed source remains
separate. The maximum 256-image source itself occupies four MiB; its boundary check
uses the normal memory layout. Smaller cases also check high addresses in both
byte orders.

Prepare an observation with \texttt{prepare.py s84-1 --source file-work.asm --output
observe.asm}. The helper changes packed source data, native controls and the output
path, leaving your instructions intact. Existing output files are refused by the
host preparation tools. This protection concerns authoring files; it does not
change the truncate-write behavior of the simulated program.

After encoding, implement \texttt{video\_store}: path in \texttt{\$a0}, admitted
byte source in \texttt{\$a1}, exact count in \texttt{\$a2} and an eight-byte-aligned
16-byte result context in \texttt{\$a3}. Admit counts one through 7\,344\,152.
The byte source also has one readable byte immediately after its exact extent:
MARS service fifteen currently reads a lookahead byte beyond each requested
window, then passes only the requested count to the file environment. The guard
is not transferred and must remain unchanged. The path is readable and
NUL-terminated; source, path, result context and private
stack are disjoint. On success return zero/count and store count at context offset
zero. On failure return one/zero and retain the context. Offset eight is always
retained. The byte source and path stay immutable. Each request transfers at most
4096 bytes, and progress follows the actual positive result. Close exactly the
descriptor acquired by this call on every completed path that owns one.

This transport routine does not validate a container and does not provide atomic
file replacement. A successful truncate-write open can discard the old target
before the first byte is written. A later error can leave an empty file or a
transferred prefix. The provided caller therefore validates and seeks through the
staged container before opening its path. It then stores and reloads the file.
Reload count equality is one transport check; the next playback stage must validate
the newly loaded bytes before trusting a directory.
